% Short chapter draft based on the local implementation, September 2026.
% Requires amsmath, booktabs, listings, and hyperref in the thesis preamble.
% Compile preview.tex from this directory for a standalone preview.
\chapter{Kappa Motion Planner Software Framework}
\label{ch:kappa-software}

\section{Overview}

The planning methods presented in the preceding chapters are implemented in
\emph{Kappa Motion Planner}, a Python package developed within the ARENA research
project. The package brings together the geometric constructions, vehicle models,
and motion primitives used by the analytical planners. Its purpose is to make
these components available through a common software interface, allowing the
same implementation to be used for numerical experiments, visualization, and
integration into a larger navigation system.

The package is imported under the name \texttt{kappa\_planner}. Its source code
is organized around five main modules: \texttt{geometry}, \texttt{corridor},
\texttt{vehicle}, \texttt{trajectory}, and \texttt{motion\_planner}. Supporting
functions are collected in the \texttt{helpers} subpackage, while executable
examples and research experiments are maintained separately. This organization
reflects the recurring structure of the analytical methods: geometric objects
describe the planning problem, vehicle parameters determine the admissible
motions, and planning routines assemble elementary primitives into a trajectory.

The main corridor-planning interface receives an ordered corridor sequence and
the boundary poses. The selection of this sequence, for example from a map or a
corridor graph, belongs to the surrounding application. Similarly, the package
provides a trajectory that an external tracking controller can use, but the
analytical planning interface does not itself implement a complete navigation
system. The core library does not require ROS, so it can also be used directly
in Python scripts.

This chapter explains how the software connects the analytical methods developed
in this thesis, introduces its principal components, and illustrates their use
in a planning example. Installation instructions are provided in
Appendix~\ref{app:kappa-installation}. The mathematical derivations and the
conditions under which the methods apply are given in the preceding chapters.

\section{From Analytical Methods to Software}

The methods developed in this thesis repeatedly use the same geometric
constructions: supporting circles, tangent connections, and straight and curved
motion primitives. Collecting these constructions in a common implementation
allows the boundary-connection routines and the corridor planners to share the
same representations. A change to a geometric construction can therefore be
applied at its source, rather than independently in each experimental script.
The unicycle and bicycle planners retain their model-specific planning rules
while using common corridor and geometric objects.

Retaining an explicit sequence of primitives also provides a direct connection
between the mathematical solution and its numerical evaluation. The structure
of a trajectory can be inspected, its traversal time evaluated from the
individual motions, and its path and controls sampled for comparison with
numerical optimal-control solutions. The experiments and plotting scripts use
these objects to study candidate families and illustrate corridor transitions.
This makes the package a shared implementation basis for the methods and their
evaluation throughout the thesis.

The software contribution is thus the organization of these analytical
constructions into reusable components and planning interfaces. Its scope
remains tied to the assumptions of the implemented methods. Those assumptions
are reflected in input checks and in the choice of planning routine; the
software does not extend the mathematical guarantees to arbitrary corridor
sequences or vehicle models.

\section{Software Organization and Main Components}
\label{sec:kappa-classes}

\subsection{Geometry and corridors}

The \texttt{Point}, \texttt{Pose}, and \texttt{Circle} classes in
\texttt{geometry.py} represent the basic planar objects. A \texttt{Pose}
combines a \texttt{Point} with a heading angle. Some interfaces also accept
poses as three-element lists or arrays containing $(x,y,\theta)$, as in the
example below.

A rectangular corridor is represented by \texttt{CorridorWorld}. Its defining
parameters are its width, length, center, and orientation. In the constructor,
the longitudinal dimension is called \texttt{height}, and \texttt{tilt} is the
angle of the forward corridor direction measured from the global $x$-axis.
From these parameters, the object computes its corners, boundary representation,
and the midpoints of its front and rear edges, stored as \texttt{head} and
\texttt{tail}. The \texttt{shrink} method constructs an inset corridor, while
\texttt{update} changes the defining parameters and recomputes the derived
geometry. These methods allow the planner and the visualization routines to use
the same geometric representation.

The corridor planners also use \texttt{IntermediateCircle} and
\texttt{IntermediateCirclesSequence}. An intermediate circle stores not only a
center and radius, but also information about its associated corridor corner,
turn direction, and admissible displacement. The sequence object maintains the
order of these circles along the route. These structures make the intermediate
constructions used by the analytical planner available for inspection and
plotting.

\subsection{Vehicle models}

The \texttt{vehicle.py} module provides \texttt{Unicycle} and \texttt{Bicycle},
derived from the base class \texttt{Vehicle}. These classes store the vehicle
dimensions, state, and input bounds, and provide the corresponding kinematic
dynamics. Parameters may be supplied explicitly or loaded from the YAML vehicle
libraries. They can subsequently be modified through the \texttt{update} method.

For the unicycle, the radius associated with simultaneous maximum translational
and angular speed is
\begin{equation}
  R=\left|\frac{v_{\max}}{\omega_{\max}}\right|.
\end{equation}
For the bicycle, the steering bound and wheelbase $L$ define
\begin{equation}
  R=\left|\frac{L}{\tan\delta_{\max}}\right|.
\end{equation}
Both quantities are stored under the implementation name \texttt{max\_radius}.
Their physical interpretation follows the corresponding vehicle model; in
particular, the unicycle also permits turns on the spot.

\subsection{Motion primitives and planning interface}

The \texttt{trajectory.py} module contains the motion-primitive classes.
Straight unicycle motions are represented by \texttt{LinearSegmentUnicycle}.
The classes \texttt{CurvilinearArcUnicycle} and \texttt{TurnOnTheSpot}
represent circular arcs and turns on the spot, respectively. Together, they
implement the primitives $S$, $C$, and $T$. Each stores its boundary
configurations and traversal time, together with sampled coordinates and
control profiles. Thus, the trajectory retains its primitive structure while
also providing data suitable for numerical evaluation and plotting. The module
also contains classes for backward arcs and sampled bicycle trajectories.

The principal corridor-planning interface is \texttt{MotionPlanner}. It takes
a vehicle, an ordered list of corridors, and initial and final poses. It
constructs the inset corridors, resolves the boundary poses, checks the relevant
input conditions, and prepares the intermediate geometry. Calling
\texttt{compute\_trajectory\_analytical()} returns an ordered collection of
motion primitives. The traversal time is obtained by summing their
\texttt{maneuver\_time} attributes; the computation time of the trajectory
generation call is recorded separately as \texttt{comp\_time\_analytical\_sol}.
The latter does not include the earlier planner construction.

The current implementation distinguishes two modes through the
\texttt{assumptions} argument: \texttt{standing} and \texttt{core}. These select
different geometric preparation and planning routines. The example below
explicitly selects \texttt{standing}, which is used for the journal corridor
construction, rather than relying on the default \texttt{core} mode.

The interface should be distinguished from the lower-level planning functions.
For example, \texttt{compute\_all\_pose\_to\_pose\_trajectories}, in
\texttt{helpers.pose\_to\_pose\_unicycle}, generates the $CSC$, $TCSC$, $CSCT$,
and $TCSCT$ candidates for the four combinations of turn directions. This
function currently requires the distance between the boundary positions to
exceed $4R$. Likewise, the existence of a geometric or vehicle class does not
imply support for every corresponding planning problem: single-corridor and
waypoint trajectory generation are not implemented through the current
\texttt{MotionPlanner} entry point.

\section{Example: Planning Through Two Corridors}
\label{sec:kappa-example}

The following example defines two overlapping corridors and computes a
trajectory for a circular unicycle. The first corridor is aligned with the
global $x$-axis. The second is constructed from the first corridor's head to a
point above it, giving a left-turn transition. The helper
\texttt{get\_corridor\_from\_vector} determines the corridor center and
orientation from these endpoints; \texttt{add\_height} adds longitudinal overlap.
The boundary poses are supplied explicitly in global coordinates, with headings
in radians.

\lstinputlisting[language=Python,float=htbp,caption={Analytical planning through two
corridors using Kappa Motion Planner.},label={lst:kappa-example}]{kappa_example.py}

The vehicle is loaded from the supplied model library, after which its speed
bounds are updated. With $v_{\max}=0.6\,\mathrm{m/s}$ and
$\omega_{\max}=1\,\mathrm{rad/s}$, the supporting turning-circle radius is
$R=0.6\,\mathrm{m}$. The planner then constructs the analytical trajectory
through the prescribed corridor sequence. The example prints both its traversal
time and the measured trajectory-generation time, and displays the resulting
path over the corridor geometry.

The returned primitive collection can also be used without plotting. Its timing,
position, heading, and control data are available for analysis or conversion to
an application-specific trajectory message. This small example illustrates the
role of the package within the thesis: it provides reusable implementations of
the analytical constructions, while allowing experiments and external
applications to define their own planning inputs and use the resulting motion.
